\section*{Chapter \ref{ch:fm:the-bank-markets-division} --- The Bank Markets Division}

\begin{solution}{exo:fm:the-bank-markets-division:1}
$307.6/5\,302.0=5.8\%$; the 5.0\% requirement holds $0.05\times5\,302.0=\$265.1$ billion of Tier 1.
\end{solution}

\begin{solution}{exo:fm:the-bank-markets-division:2}
$27\,761/149\,500=18.6\%$. The bank's segment ROE uses net income applicable to common equity, after the preferred dividends it allocates to the
segment, which is lower than net income.
\end{solution}

\begin{solution}{exo:fm:the-bank-markets-division:3}
After-tax profit $1.4\times0.75=\$1.05$ billion. By RWA: $1.05/(0.13\times15)=1.05/1.95=53.8\%$. By leverage: $1.05/(0.05\times400)=1.05/20=5.2\%$.
\end{solution}

\begin{solution}{exo:fm:the-bank-markets-division:4}
$0.12\times0.05=0.6\%$ a year of leverage exposure. Break-even: $1.4/400=0.35\%$ (35 basis points).
\end{solution}

\begin{solution}{exo:fm:the-bank-markets-division:5}
Desk maxima: $12.5+20.0+10.0+9.1+4.0+6.5=62.1$. The group needs $\max(35.1, 55.5)=55.5$. Equity derivatives and credit are bound by RWA while the
group is bound by leverage, so their RWA equity is counted on top of a leverage requirement that already covers it.
\end{solution}

\begin{solution}{exo:fm:the-bank-markets-division:6}
Under the binding allocation: prime services (13.5\%), equity derivatives (18.1\%) and FX (28.1\%). Under the RWA key only: flow rates (15.4\%) and
repo (53.8\%).
\end{solution}

\begin{solution}{exo:fm:the-bank-markets-division:7}
Profit falls to \$5.625 billion (from \$7.05 billion), equity to \$48.1 billion, and the return to 11.7\%: most of the gain in return from cutting
repo is lost once more of the prime business follows it.
\end{solution}

\begin{solution}{exo:fm:the-bank-markets-division:8}
The programme assumes that the other desks can grow by half at unchanged margins and that the desks are independent; clients who finance with the
bank also trade and clear with it. Cut in steps, measure the effect on prime and flow revenue, and reprice repo with a \omterm{def:fm:the-bank-markets-division:charge}{balance-sheet charge}
before shrinking it.
\end{solution}

\begin{solution}{pb:fm:the-bank-markets-division:1}
\begin{enumerate}
\item See \cref{def:fm:the-bank-markets-division:cet1,def:fm:the-bank-markets-division:lr}.
\item 4.5\% minimum, 2.5\% \omterm{def:fm:the-bank-markets-division:buffers}{stress capital buffer}, 4.5\% GSIB surcharge (method 2).
\item Risk-based Tier 1: $13.0\%\times1\,981.7=\$257.6$ billion; leverage: $5.0\%\times5\,302.0=\$265.1$ billion; leverage binds, narrowly.
\item The static enhanced-SLR buffers were replaced by 50\% of the holding company's method 1 GSIB surcharge, capped at 1\% for bank subsidiaries.
\item By Standardized RWA and the GSIB surcharge, and by a simulation of capital depletion in a severe stress.
\item See \cref{def:fm:the-bank-markets-division:alloc}.
\item 53.8\%, 5.2\% and 5.2\%.
\item Equity derivatives and credit: their risk weights are high relative to their balance sheet.
\item See \cref{prop:fm:the-bank-markets-division:max}.
\item \$35.1, 55.5 and 62.1 billion; 11.4\% on \$62.1 billion.
\item The cost of the equity the leverage rule holds per unit of leverage exposure: $0.12\times0.05=0.6\%$ a year.
\item Flow rates 0.075, repo $-0.75$, prime services 0.45, equity derivatives 1.11, FX 0.765, credit 0.405 (\$ billion).
\item 0.35\% of leverage exposure a year.
\item It shrinks repo to about 0.11 of its size and grows the other desks to 1.5 times; profit \$9.12 billion against \$7.05 billion.
\item That desks can grow without lowering margins, and that they are independent.
\item Profit \$6.08 billion (from 7.05), equity \$50.1 billion (from 62.1), return 12.1\% (from 11.4\%).
\item Their binding constraints differ, so the same business costs them different amounts of equity.
\item How much of its revenue comes from clients who also finance with the repo desk, and how much of it would move with the financing.
\item 53.8\% by RWA, 5.2\% by leverage; break-even \omterm{def:fm:the-bank-markets-division:charge}{balance-sheet charge} 35 basis points of leverage exposure a year.
\item It uses little risk capital and a great deal of balance sheet, and for this division the balance sheet is what binds. Priced at the cost of
the equity the leverage rule holds, it does not earn its hurdle.
\end{enumerate}
\end{solution}

\begin{solution}{iq:fm:the-bank-markets-division:1}
A risk-based ratio weights exposures by risk; a \omterm{def:fm:the-bank-markets-division:lr}{leverage ratio} does not. The \omterm{def:fm:the-bank-markets-division:lr}{leverage ratio} is a backstop against risk weights that are too low or
gamed.
\iqlookfor{the backstop role and that each binds for different businesses.}
\end{solution}

\begin{solution}{iq:fm:the-bank-markets-division:2}
Repo assets are short, secured loans with low risk weights, but they are on the balance sheet at full size.
\iqlookfor{risk weights against balance-sheet size.}
\end{solution}

\begin{solution}{iq:fm:the-bank-markets-division:3}
Equity held $0.05\times200=\$10$ billion; cost at 12\% is \$1.2 billion a year, more than the \$1 billion it earns: it destroys value.
\iqlookfor{equity from the leverage rule times the hurdle.}
\end{solution}

\begin{solution}{iq:fm:the-bank-markets-division:4}
The sum of desk-level maxima exceeds the maximum of the sums unless one constraint binds for every desk
(\cref{prop:fm:the-bank-markets-division:max}).
\iqlookfor{$\sum\max\ge\max\sum$.}
\end{solution}

\begin{solution}{iq:fm:the-bank-markets-division:5}
Different balance sheets make different constraints bind: a leverage-bound bank prices balance sheet high and risk low, and the reverse.
\iqlookfor{the binding constraint as a competitive position.}
\end{solution}

\begin{solution}{iq:fm:the-bank-markets-division:6}
Margins that fall as a desk grows, links between desks (clients who finance and trade), the stress constraint, limits on how fast desks can grow,
and the cost of shrinking a franchise.
\iqlookfor{concave revenue, franchise links and adjustment costs.}
\end{solution}
